% ==========================================================================
% GENERATED by scripts/split_paper.py. Do not edit this file.
%
% Source: imported_papers/grokking/iclr2027_conference_v2.tex
% Part 14: Supplementary analysis of generalization
%
% Edits belong upstream, in the published paper. To change what the
% chapter includes, change imports/grokking.toml and re-run
% the splitter.
% ==========================================================================
\section{Supplementary analysis of generalization}
\label{app:generalization}

Section~\ref{sec:grokrules} reports that a weight-space measurement, $q$, separates the runs that generalize from those that only fit the training set. This appendix gives the definition, the compared configurations, the full results and the limits of the measurement.

\paragraph{A local measure on the weight vector.} A well-studied account of grokking explains generalization through weight decay: a generalizing circuit fits the training data with smaller weights than a memorizing one, so regularization gradually favors it \citep{nanda2023progress,varma2023circuit,liu2022omnigrok}. Once the training data are fitted, the network can keep moving among parameter settings that preserve the fit, and weight decay pushes it toward the smallest of them, unless some other component of the learning signal opposes it. To measure this, let $J(\theta)$ be the Jacobian of all training logits with respect to the parameters. A small parameter change $\vv$ changes the logits by $J\vv$ to first order, so the directions with $J\vv = 0$, the null space of $J$, are the directions in which weights can be reduced without changing the training outputs to first order. Let $P$ be the orthogonal projector onto this null space. Then $P\theta$ is the component of the current weight vector along these directions, and we measure its norm,
\begin{equation}
q(\theta) \;=\; \|P\theta\|_2 .
\label{eq:flatmass}
\end{equation}
Both $J$ and $P$ depend on the current parameters, so $q$ is a local measurement that we recompute at every checkpoint; it is not a count of unnecessary parameters or neurons.

For continuous dynamics $\dot{\theta} = -\nabla \ell - \lambda\theta - \vb$, where $\ell$ is the training loss, $\lambda > 0$ the weight-decay strength and $\vb$ the learning signal's departure from the loss gradient, the loss gradient has no component in the null space and
\begin{equation*}
P\dot{\theta} \;=\; -\lambda P\theta - P\vb .
\end{equation*}
Weight decay shrinks the projected weights while an additional signal can oppose it. If the Jacobian is held fixed and $P\vb = 0$, the projected weights decay exponentially and $q$ tends to zero; a constant nonzero $P\vb$ instead leaves the stationary value $P\theta = -P\vb/\lambda$ (Proposition~\ref{prop:persist}). This calculation motivates measuring $q$ during training. Since the projector changes and Adam transforms the signal, we measure $q$ for every rule, including backpropagation.

\paragraph{Compared configurations.} We measure nine configurations on modular addition at $p = 31$ with three seeds each and a $15{,}000$-epoch budget: the reference, full credit equalization, DFC, feedback alignment (FA) and direct feedback alignment (DFA) \citep{lillicrap2016feedback,nokland2016dfa}, sign-SGD \citep{bernstein2018signsgd}, and three constructed controls. Every configuration writes its learning signal into the gradient slot, and the signal is then clipped and passed to Adam at the shared learning rate $10^{-3}$ and weight decay $10^{-3}$. Sign-SGD also memorizes when applied directly, at learning rates $10^{-4}$ and $10^{-5}$ and with weight decay up to $10^{-2}$ (three seeds, test accuracy $0.000$). Each control replaces the update by the unit vector $\cos\varphi\,\hat{\vg} + \sin\varphi\,\hat{\vv}_\perp$ rescaled to the gradient norm, where $\hat{\vg}$ is the normalized cross-entropy gradient, $\hat{\vv}_\perp$ is a random vector projected orthogonally to the current gradient and normalized, and $\cos\varphi = 0.4$. \emph{Fresh noise} draws the random vector anew at every step. \emph{Fixed bias} draws it once per run, so the added component keeps its direction up to the projection. \emph{Error-gated bias} uses the same fixed vector but multiplies the orthogonal component by the current loss relative to its initial value before the renormalization, so its cosine to the gradient is $0.4$ at the start and rises toward one as the loss falls. Fresh noise and error-gated bias grok, fixed bias memorizes. The recorded checkpoints, at epochs $75$, $150$, $300$, $450$, $750$ and $1500$, show the delay in every grokking run: training accuracy is $1.0$ by epoch $150$ (by epoch $300$ for fresh noise), and test accuracy first exceeds $0.95$ between epochs $750$ and $1500$ for the reference and error-gated bias, after epoch $1500$ for credit equalization and fresh noise, and between epochs $450$ and $750$ for DFC, whose delays at $10$-epoch resolution are $520$ to $610$ epochs (Appendix~\ref{app:repro}).

\paragraph{Results.} At epoch $300$, $2\%$ of training, all $27$ runs have training accuracy $1.0$ and test accuracy $0.0$. The estimated $q$ nevertheless already separates the eventual outcomes: the largest value among the runs that later grok is $22.02$, below the smallest value among the runs that memorize, $27.71$ (Table~\ref{tab:early_flatmass}). The separation is also present at the endpoints, where every grokking run has a smaller $q$ than every memorizing run (Table~\ref{tab:flatmass}). Total parameter norm does not separate the runs at epoch $300$: the largest norm among the runs that later grok, $42.2$, exceeds the smallest norm among the memorizing runs, $41.3$. At the endpoints the norm separates the outcomes as well (grokking runs between $13.1$ and $26.7$, memorizing runs between $35.6$ and $150$). These observations agree with the weight-decay account.

\paragraph{Within one rule.} The nine configurations differ in their rule, and each one groks in every seed or memorizes in every seed, so a separation across them could reflect the rule rather than the outcome. We therefore vary one setting within two rules, with five seeds per setting and the budget, checkpoints and estimator above (Table~\ref{tab:withinrule}). For backpropagation with Adam over weight decays from $0$ to $3\times10^{-3}$, every run with weight decay groks and the runs without it do not generalize within the budget; at weight decay $10^{-2}$ the runs do not fit the training set and are excluded. $q$ at epoch $300$ is lower in every grokking run than in every run without decay, but the runs at weight decay $10^{-4}$ grok with $q = 31.3$, above the smallest value among the memorizing configurations of Table~\ref{tab:early_flatmass}. For the fixed-bias rule, raising the cosine between update and gradient from $0.4$ to $0.95$ lowers $q$ at epoch $300$ from $48.8$ to $18.6$ and raises final test accuracy from $0.00$ to $0.98$, with a Spearman correlation of $-0.97$ over the $40$ runs. Runs at cosines up to $0.8$ do not meet the grokking criterion, although test accuracy reaches $0.46$ to $0.89$ at $0.7$ and $0.8$, and at $0.85$, $0.9$ and $0.95$ two, four and four of five seeds grok. No single threshold on $q$ at epoch $300$ separates the runs that grok from the others in this family (largest value among the grokking runs $24.9$, smallest among the others $18.6$), and within a setting the differences in $q$ between seeds, about $0.1$, do not consistently order the outcomes. Within a rule, $q$ thus changes with a setting in the direction the weight-decay account predicts, but a common threshold on $q$ early in training does not predict the outcome.

\begin{table}[!ht]
\caption{Null-space weight norm $q$ within one rule (modular addition, $p=31$, five seeds per setting, $15{,}000$ epochs; standard deviations in parentheses). Grokking seeds meet the criterion of Section~\ref{sec:setup}; final test is the mean test accuracy at the end of training.}
\label{tab:withinrule}
\begin{center}
\small
\begin{tabular}{lcccc}
\toprule
Setting & $q$ at epoch $300$ & $q$ at epoch $750$ & Grokking seeds & Final test \\
\midrule
\multicolumn{5}{l}{\emph{Backpropagation with Adam, weight decay}} \\
$0$ & $35.97\,(0.13)$ & $37.33$ & $0/5$ & $0.000$ \\
$10^{-4}$ & $31.27\,(0.07)$ & $22.91$ & $5/5$ & $0.995$ \\
$3\times10^{-4}$ & $24.96\,(0.05)$ & $13.44$ & $5/5$ & $0.997$ \\
$10^{-3}$ & $14.34\,(0.08)$ & $7.57$ & $5/5$ & $0.991$ \\
$3\times10^{-3}$ & $8.38\,(0.10)$ & $4.96$ & $5/5$ & $0.996$ \\
$10^{-2}$ & $5.71\,(0.12)$ & $1.69$ & training set not fitted & $0.162$ \\
\midrule
\multicolumn{5}{l}{\emph{Fixed-bias rule, cosine between update and gradient}} \\
$0.4$ & $48.82\,(0.19)$ & $49.71$ & $0/5$ & $0.000$ \\
$0.5$ & $42.80\,(0.15)$ & $42.21$ & $0/5$ & $0.000$ \\
$0.6$ & $37.50\,(0.10)$ & $35.42$ & $0/5$ & $0.026$ \\
$0.7$ & $32.58\,(0.12)$ & $29.08$ & $0/5$ & $0.548$ \\
$0.8$ & $27.56\,(0.08)$ & $22.81$ & $0/5$ & $0.837$ \\
$0.85$ & $24.91\,(0.06)$ & $19.58$ & $2/5$ & $0.941$ \\
$0.9$ & $21.99\,(0.04)$ & $16.27$ & $4/5$ & $0.972$ \\
$0.95$ & $18.64\,(0.06)$ & $12.47$ & $4/5$ & $0.981$ \\
\bottomrule
\end{tabular}
\end{center}
\end{table}

\begin{table}[!ht]
\caption{Estimated null-space weight norm $q$ early in training, as means (standard deviations) over three seeds. Every run fits the training set and has zero test accuracy at epoch $300$. The outcome column refers to the completed runs.}
\label{tab:early_flatmass}
\begin{center}
\small
\begin{tabular}{lccc}
\toprule
Configuration & Epoch $300$ ($2\%$) & Epoch $750$ ($5\%$) & Outcome \\
\midrule
Reference (backpropagation with Adam) & $14.32\,(0.09)$ & $7.63\,(0.14)$ & groks \\
Credit equalization ($\beta{=}1$) & $14.46\,(0.08)$ & $9.11\,(0.13)$ & groks \\
DFC & $11.69\,(0.14)$ & $6.48\,(0.08)$ & groks \\
Error-gated bias & $21.93\,(0.09)$ & $9.53\,(0.04)$ & groks \\
Fresh noise & $15.37\,(0.26)$ & $9.32\,(0.13)$ & groks \\
DFA & $27.80\,(0.10)$ & $26.05\,(0.36)$ & memorizes \\
FA & $28.45\,(0.21)$ & $27.63\,(0.31)$ & memorizes \\
Fixed bias & $48.94\,(0.08)$ & $49.71\,(0.05)$ & memorizes \\
Sign-SGD & $49.13\,(0.02)$ & $72.33\,(0.48)$ & memorizes \\
\bottomrule
\end{tabular}
\end{center}
\end{table}

\begin{table}[!ht]
\caption{Estimated null-space weight norm $q$ at the end of training ($15{,}000$ epochs of modular addition at $p=31$; means over three seeds). Every grokking configuration has a smaller $q$ than every memorizing one. At initialization, $q = 29.9$.}
\label{tab:flatmass}
\begin{center}
\small
\begin{tabular}{lcllcl}
\toprule
Configuration & $q(\theta)$ & Outcome & Configuration & $q(\theta)$ & Outcome \\
\midrule
Fresh noise & $5.30$ & groks & DFA & $26.34$ & memorizes \\
Reference & $5.90$ & groks & FA & $26.55$ & memorizes \\
Error-gated bias & $6.55$ & groks & Fixed bias & $51.80$ & memorizes \\
Credit equalization ($\beta{=}1$) & $7.24$ & groks & Sign-SGD & $134.72$ & memorizes \\
DFC & $8.49$ & groks & & & \\
\bottomrule
\end{tabular}
\end{center}
\end{table}

\paragraph{Numerical estimation.} We estimate the projection as
\[
\widehat{P\theta} \;=\; \theta - J^\top (J J^\top + \epsilon I)^{-1} J\theta ,
\]
where the subtracted term approximates the component of $\theta$ in the span of the training-logit gradients and the damping $\epsilon > 0$ makes the remainder an approximation rather than an exact projection. The linear system is solved with conjugate gradients using Jacobian-vector products, without forming $J$. The reported measurements use relative damping $10^{-6}$ and $400$ conjugate-gradient iterations. On small networks the estimate agrees with an explicit singular-value decomposition to a relative error below $10^{-4}$. A longer solve at one trained width-$256$ reference checkpoint ($10{,}000$ preconditioned iterations) gives $q = 6.04$ with a relative residual of $10^{-3}$ and an idempotence error of $0.010$ relative to the parameter norm. At epoch $300$, retraining the nine configurations and re-estimating $q$ with relative damping from $10^{-10}$ to $10^{-4}$, up to $10{,}000$ iterations and in double precision changes every value by at most $0.01$, and the largest value among the runs that later grok stays below the smallest among those that memorize under every setting.

\begin{proposition}[Projected weight dynamics with a fixed local Jacobian]
\label{prop:persist}
Let $\ell$ be a differentiable loss of the training logits $z(\theta)$, let $J = \partial z / \partial \theta$, and let $P$ project orthogonally onto $\operatorname{null}(J)$. Under $\dot{\theta} = -\nabla_\theta \ell - \lambda\theta - \vb$ with $\lambda > 0$,
\begin{equation}
P\dot{\theta} \;=\; -\lambda P\theta - P\vb .
\label{eq:null_dynamics}
\end{equation}
At a stationary point, $P\theta = -P\vb/\lambda$. In the local model in which $J$ is held fixed, write $\vv = P\theta$. If $P\vb = \vb_0$ is constant, then
\[
\vv(t) \;=\; e^{-\lambda t}\vv(0) - \frac{1 - e^{-\lambda t}}{\lambda}\,\vb_0 ,
\]
so $\|\vv(t)\|_2$ tends to $\|\vb_0\|_2/\lambda$. If $P\vb(t)$ is locally integrable and tends to zero, then $\vv(t)$ tends to zero.
\end{proposition}
\begin{proof}
The chain rule gives $\nabla_\theta \ell = J^\top \nabla_z \ell$, so $P \nabla_\theta \ell = 0$. Projecting the dynamics gives Equation~\ref{eq:null_dynamics} and the stationary-point condition. For fixed $J$, the integrating-factor solution is
\[
\vv(t) \;=\; e^{-\lambda t}\vv(0) - \int_0^t e^{-\lambda (t-s)} P\vb(s)\, ds .
\]
Constant forcing gives the displayed expression. If the forcing tends to zero, splitting the integral into a fixed initial interval and the remaining tail shows that both contributions tend to zero.
\end{proof}
The proposition concerns a continuous-time model with a fixed Jacobian and the signal applied directly.
